% !PRINTED 669
\weberpage{669}
%<<CONT
and by multiplication of (6) and (7), with regard to (4)

\weq{8}{e(r^n)\, e(r)^{-n} = \mathfrak{E}(r)^2.}

In our considerations on the class number the following theorem will yet appear:

\gesperrt{A unit $e(r)$ of the real field $H_m$, which is positive together with all its conjugates, is the square of a unit in the field $H_m$.}

But formula (8) shows that the unit $\pm e(r)$ has this property, and hence that $e(r)$ (since $-1 = i^2$ is likewise the square of a unit) is the square of a unit of the field $\Omega_m$.

With this the question is reduced to the proof of the two theorems mentioned. That proof appears as the keystone of a long chain of considerations, extremely rich in interesting relations, to which the next chapters are to be devoted.

\begin{center}
\rule{6em}{0.4pt}
\end{center}
